% ECONOMICS_PROBLEM: {"id": "C72-589", "title": "Average-case analysis of specified search-and-mix algorithms", "jel_code": "C72", "source_id": "OP-0134"}
% !TeX program = xelatex
\documentclass[11pt,a4paper]{article}
\usepackage[margin=23mm]{geometry}
\usepackage{fontspec}
\setmainfont{Latin Modern Roman}
\usepackage{amsmath,amssymb,mathtools}
\usepackage{xcolor,enumitem,booktabs,longtable,array,xurl,hyperref}
\definecolor{muted}{HTML}{53636C}
\hypersetup{colorlinks=true,urlcolor=blue,linkcolor=blue}
\setlength{\parindent}{0pt}
\setlength{\parskip}{5pt}
\newcommand{\R}{\mathbb R}
\newcommand{\E}{\mathbb E}
\newcommand{\N}{\mathbb N}
\newcommand{\ind}{\mathbf 1}
\newcommand{\Var}{\operatorname{Var}}
\newcommand{\Cov}{\operatorname{Cov}}
\newcommand{\supp}{\operatorname{supp}}
\DeclareMathOperator*{\argmin}{arg\,min}
\DeclareMathOperator*{\argmax}{arg\,max}
\newcommand{\entrymeta}[3]{{\sffamily\small\color{muted}\textbf{\texttt{#1}}\quad|\quad #2\par #3\par}}
\newcommand{\Problemref}[1]{\texttt{#1}}
\newcommand{\JELref}[2]{\texttt{#2} (JEL \texttt{#1})}
\begin{document}
\section*{Average-case analysis of specified search-and-mix algorithms}
\textbf{Problem ID:} \texttt{C72-589}\quad \textbf{JEL:} \texttt{C72}\par
% BEGIN ECONOMICS STATEMENT
\label{problem:OP-0134}
% GAME_THEORY_PROBLEM: {"local_id": "GT-004", "original_number": 4, "kind": "R", "entry_kind": "statement_only", "published": false, "review_required": true, "category": "game_theory"}
\label{game:GT-004}
{\sffamily\small\color{muted}
\textbf{GT-004} | Algorithmic equilibrium and complexity\\
\textbf{R} | Parameterized average-case research target.
\par}
\subsection*{Definitions and mathematical statement}
Specify a sequence $(\mathcal D_n)_{n\geq1}$ of distributions on rational $n\times n$ bimatrix games with payoffs in $[0,1]$, a sampling and bit-encoding convention, and a particular search-and-mix algorithm $\mathcal A$, including its tolerances and internal randomization. Let $\xi$ denote the algorithm's random seed, independent of the sampled game. For its rational output $Z_n=\mathcal A(G,\xi)$, let $r_G(Z_n)$ be the largest unilateral payoff gain and $T_{\mathcal A}(G,\xi)$ its bit-operation count. Define
\[
 e_n=\mathbb E_{G\sim\mathcal D_n,\xi}[r_G(Z_n)],\qquad
 \tau_n=\mathbb E_{G\sim\mathcal D_n,\xi}[T_{\mathcal A}(G,\xi)].
\]
For each fully specified pair $(\mathcal D,\mathcal A)$, determine asymptotic equivalents, or matching upper and lower asymptotic orders, for $e_n$ and $\tau_n$. A finite-encoding benchmark is independent uniform entries in $\{0,2^{-b(n)},\ldots,1\}$ for a declared integer function $b(n)$, using a declared initialization law and tolerance sequence.
\subsection*{Short English statement}
Explain rigorously how accurate and how fast a particular Nash algorithm is on a particular random-game distribution.
\subsection*{Variants and scope boundaries}
Both the distribution and the algorithm are essential data. Existence of some unspecified ``natural'' distribution would be vacuous. Uniformly bounded worst-case residuals immediately give bounded $e_n$; the substantive target is sharp distribution-sensitive behavior. Arithmetic-operation counts and bit-operation counts need not have identical asymptotics.
\subsection*{Source relationship and status}
[S1] proposes average-case approximation analysis without selecting a universal distribution or asymptotic law. The displayed expectations are editorial definitions; the finite grid is an explicit optional benchmark rather than a model asserted to occur in the source.
\subsection*{References}
\begingroup\footnotesize\raggedright
{\footnotesize\raggedright
\textbf{[S1]} Hanyu Li, Wenhan Huang, Zhijian Duan, David Henry Mguni, Kun Shao, Jun Wang, and Xiaotie Deng (2023). \emph{A survey on algorithms for Nash equilibria in finite normal-form games}. Locator: concluding ``Open problems and future directions,'' theoretical item 3; Section 4 for experimental distributions. \url{https://arxiv.org/pdf/2312.11063}\par}
\par\endgroup

\subsection*{Common conventions: Source mathematical conventions}

\textbf{Finite games.} For a finite set $A$, $\Delta(A)$ is its probability
simplex. Players choose independent mixed strategies in Nash equilibrium;
correlation is allowed only when explicitly part of the solution concept.
In a game with payoff functions $u_i$ and strategy profile $x$, an additive
$\varepsilon$ Nash equilibrium satisfies
\[
 u_i(x)\geq u_i(x_i',x_{-i})-\varepsilon
 \quad\text{for every player $i$ and every admissible deviation $x_i'$}.
\]
For numerical approximation, payoffs are normalized as locally specified.
An $\varepsilon$ payoff residual is distinct from distance at most
$\varepsilon$ to the set of exact equilibria. Point convergence, convergence
of empirical play, and convergence in distance to an equilibrium set are
also distinct.

\textbf{Computational representations.} Unless stated otherwise, a complexity
question takes rational data with binary encoding; polynomial time is measured
in total bit length. A fixed-accuracy polynomial algorithm is not necessarily
polynomial in $1/\varepsilon$, and neither is necessarily polynomial in
$\log(1/\varepsilon)$. Succinct games and value, demand, or sampling oracles
must declare their interfaces. Exact real solutions require an explicit
representation; an unspecified list of arbitrary real numbers is not a
finite computational output. A claimed algorithmic barrier is meaningful
only for its specified model and reductions.

\textbf{Dynamic games.} Histories, information, observation, and allowable
randomization are part of the model. A stationary strategy depends only on
the current state; a positional strategy in a deterministic graph game
chooses one action at each controlled vertex. Nash equilibrium at an initial
state and equilibrium after every history are different requirements.
An expectation of a pathwise $\liminf$ payoff, a $\liminf$ of expected averages,
a limiting expected average, and a uniform finite-horizon equilibrium are
different criteria. None is silently substituted for another.

\textbf{Allocation and mechanisms.} A complete allocation assigns every item
exactly once unless otherwise stated. Zero-valued goods, negative values,
capacity constraints, feasibility, lotteries, and copies of goods affect
fairness definitions and are specified locally. Dominant-strategy incentive
compatibility, Bayesian incentive compatibility, universal truthfulness,
and truthfulness in expectation impose different inequalities. Whenever
an objective ratio has zero denominator, its convention is stated or those
instances are excluded.

\textbf{Infinite-dimensional models.} Expectations, integrals, stochastic
equations, and optimization objectives require their local regularity,
integrability, filtrations, and admissible domains. A backward--forward
mean-field system is a boundary-value problem; it is not an initial-value
semiflow without an additional construction. Long-time, large-population,
and vanishing-noise limits require an order of limits and a convergence
topology. If a minimum need not be attained, the objective is its infimum
or supremum and the approximation requirement is stated separately.
% END ECONOMICS STATEMENT
\end{document}
